\documentclass[11pt, a4paper]{article}

% ── pacchetti ─────────────────────────────────────────────────────────────────
\usepackage[T1]{fontenc}
\usepackage[utf8]{inputenc}
\usepackage[italian]{babel}
\usepackage{lmodern}
\usepackage{microtype}
\usepackage{geometry}
\geometry{margin=2.5cm, top=3cm, bottom=3cm}

\usepackage{amsmath, amssymb, bm}
\usepackage{booktabs}
\usepackage{array}
\usepackage{multirow}
\usepackage{xcolor}
\usepackage{hyperref}
\usepackage{tikz}
\usetikzlibrary{arrows.meta, positioning, shapes.geometric, fit, backgrounds}

% ── colori ────────────────────────────────────────────────────────────────────
\definecolor{myblue}{HTML}{2C5F8A}
\definecolor{mygreen}{HTML}{2E7D5B}
\definecolor{myorange}{HTML}{C25B00}
\definecolor{mygray}{HTML}{5A5A5A}
\definecolor{mylightblue}{HTML}{D9E8F5}
\definecolor{mylightgreen}{HTML}{D5EDE3}
\definecolor{mylightorange}{HTML}{FAE6D0}

% ── hyperref ──────────────────────────────────────────────────────────────────
\hypersetup{
  colorlinks = true,
  linkcolor  = myblue,
  citecolor  = myblue,
  urlcolor   = myblue
}

% ── stile sezioni ─────────────────────────────────────────────────────────────
\usepackage{titlesec}
\titleformat{\section}{\large\bfseries\color{myblue}}{\thesection.}{0.5em}{}[\vspace{-2pt}\rule{\linewidth}{0.4pt}]
\titleformat{\subsection}{\normalsize\bfseries\color{mygray}}{\thesubsection.}{0.5em}{}

% ── comandi matematici ────────────────────────────────────────────────────────
\newcommand{\N}{\mathcal{N}}
\newcommand{\HN}{\mathcal{N}^+}
\newcommand{\E}{\mathbb{E}}
\newcommand{\Var}{\mathrm{Var}}
\newcommand{\Cov}{\mathrm{Cov}}
\newcommand{\SOC}{\mathrm{SOC}}
\renewcommand{\vec}[1]{\bm{#1}}

% ── intestazione ──────────────────────────────────────────────────────────────
\title{%
  {\large Descrizione del modello}\\[4pt]
  {\LARGE\color{myblue} M-SP: Slope Proporzionale}\\[4pt]
  {\normalsize Analisi bayesiana dei profili verticali di SOC, N e P}
}
\author{}
\date{\today}

% =============================================================================
\begin{document}
\maketitle
\thispagestyle{empty}
\tableofcontents
\bigskip

% =============================================================================
\section{Motivazione}

Il carbonio organico del suolo (SOC), l'azoto totale (N) e il fosforo totale (P)
decadono tipicamente con la profondità seguendo una legge di potenza:
$y \propto \text{Bottom}^{\beta}$, linearizzabile con la trasformazione log-log
($\log y \sim \beta \log\text{Bottom}$).

Un'analisi esplorativa (OLS per campo, BLUP lmer) rivela che:
\begin{itemize}
  \item la correlazione tra l'intercetta per campo (livello a profondità media)
        e la slope (velocità di decadimento) è fortemente positiva per il SOC
        ($r \approx {+}0.93$ con lmer, Bottom centrato);
  \item i campi con SOC alto alla profondità media decadono \emph{più lentamente};
  \item per N questa correlazione è trascurabile; per P è debole ma presente.
\end{itemize}

Il modello M-SP cattura direttamente questo pattern: la deviazione del campo
dalla media globale ha un'ampiezza che cambia con la profondità in modo
proporzionale al livello del campo.

% =============================================================================
\section{Dati}

\subsection{Struttura}

$N \approx 220$ osservazioni, $J = 40$ campi, fino a 6 profondità per campo.

\medskip
\begin{tabular}{@{}lll@{}}
\toprule
File & Formato & Contenuto \\
\midrule
\texttt{data/crop.csv} & CSV & dati grezzi \\
\texttt{data/dati.rds} & R & preprocessati (ILR, log, scala) \\
\bottomrule
\end{tabular}

\subsection{Variabili risposta}

\[
  y_{r,ij} = \log(\text{risposta}_{r,ij}), \quad r \in \{\SOC,\, N,\, P\}
\]

\subsection{Covariati within-field $X_W$ ($N \times 4$)}

\medskip
\begin{tabular}{@{}clll@{}}
\toprule
$k$ & Covariato & Descrizione & Trasformazione \\
\midrule
1 & \texttt{logBottom} & $\log(\text{profondità [cm]})$ & standardizzato \\
2 & \texttt{Texture1}  & ILR\textsubscript{1}(Clay, Silt, Sand) & standardizzato \\
3 & \texttt{Texture2}  & ILR\textsubscript{2}(Clay, Silt, Sand) & standardizzato \\
4 & \texttt{BulkDensity} & densità apparente & standardizzato \\
\bottomrule
\end{tabular}

\medskip
Le coordinate ILR (Isometric Log-Ratio) eliminano il vincolo compositionale
$\text{Clay} + \text{Silt} + \text{Sand} = 100\%$. Calcolate via
\texttt{compositions::ilr(acomp(Clay, Silt, Sand))} con base di Helmert standard:
\[
  T_1 = \sqrt{\tfrac{1}{2}} \log\!\frac{\text{Clay}}{\text{Silt}},
  \qquad
  T_2 = \sqrt{\tfrac{2}{3}} \log\!\frac{\sqrt{\text{Clay}\cdot\text{Silt}}}{\text{Sand}}.
\]

\subsection{Covariati between-field $X_B$ ($J \times 4$)}

\medskip
\begin{tabular}{@{}cll@{}}
\toprule
$k$ & Covariato & Descrizione \\
\midrule
1 & \texttt{OnFarm} & gestione organica in loco (0/1) \\
2 & \texttt{Irrigate} & irrigazione (0/1) \\
3 & \texttt{Fertilised} & fertilizzazione minerale (0/1) \\
4 & \texttt{N\_Natural} & apporto naturale di azoto (0/1) \\
\bottomrule
\end{tabular}

% =============================================================================
\section{Struttura del modello}

\subsection{Likelihood}

Per ogni risposta $r$, osservazione $i$ nel campo $j$:
\[
  y_{r,ij} \sim \N\!\left(\mu_{r,ij},\; \sigma_r^2\right)
\]

\subsection{Predittore lineare}

\[
  \boxed{
  \mu_{r,ij} =
    \underbrace{\alpha_r}_{\text{media globale}}
    +\;
    \underbrace{z_{\nu,r,j}\,\bigl(\psi_r + \eta_r \cdot x_{\text{logB},i}\bigr)}_{\text{RI + slope proporzionale}}
    +\;
    \underbrace{\vec{\gamma}_r^\top \vec{x}_{W,i}}_{\text{within fissi}}
    +\;
    \underbrace{\vec{\beta}_r^\top \vec{x}_{B,j}}_{\text{between fissi}}
  }
\]

dove $x_{\text{logB},i}$ è la prima colonna di $X_W$ (logBottom standardizzato).

\subsection{Parametrizzazione non centrata (NCP)}

\[
  z_{\nu,r,j} \sim \N(0,\,1), \qquad j = 1,\ldots,J
\]

L'effetto casuale vero è $\nu_{r,j} = z_{\nu,r,j}\,(\psi_r + \eta_r \cdot x_{\text{logB},i})$,
ma a profondità media ($x_{\text{logB}} = 0$):
$\nu_{r,j}\big|_{x=0} = z_{\nu,r,j}\,\psi_r$.

La NCP è necessaria per l'efficienza del campionatore NUTS con prior diffusi
sugli effetti casuali.

\subsection{Interpretazione di $\eta_r$}

\[
  \Var[\nu_{r,j} \mid x_{\text{logB}}]
  = \left(\psi_r + \eta_r \, x_{\text{logB}}\right)^2
\]

\begin{itemize}
  \item $\eta_r > 0$: la variabilità between-field \emph{aumenta} alle profondità
        maggiori → i campi più ricchi a profondità media decadono più lentamente.
  \item $\eta_r = 0$: profilo verticale uguale per tutti i campi (slope fissa).
  \item Il coefficiente di proporzionalità è $b_r = \eta_r / \psi_r$.
\end{itemize}

\subsection{Intercetta globale $\alpha_r$}

L'effetto casuale $z_{\nu,r,j}$ ha media zero per costruzione (NCP).
Senza $\alpha_r$, i coefficienti $\vec{\beta}_r$ (management) assorbirebbero
la media globale di $\log r$, esplodendo a valori $|\beta| \approx 2$.
L'intercetta globale è quindi un componente \emph{necessario} del modello.

% =============================================================================
\section{Modello completo}

\subsection{Parametri e prior}

\medskip
\begin{tabular}{@{}lll@{}}
\toprule
Parametro & Dimensione & Prior \\
\midrule
$\alpha_r$ & scalare & $\N(0,\;1)$ \\
$z_{\nu,r,j}$ & vettore $J$ & $\N(0,\;1)$ (NCP) \\
$\psi_r$ & scalare $>0$ & $\HN(0,\;1)$ \\
$\eta_r$ & scalare & $\N(0,\;1)$ \\
$\vec{\gamma}_r$ & vettore $K_W = 4$ & $\N(\vec{0},\;\vec{I})$ \\
$\vec{\beta}_r$ & vettore $K_B = 4$ & $\N(\vec{0},\;\vec{I})$ \\
$\sigma_r$ & scalare $>0$ & $\HN(0,\;1)$ \\
\midrule
Totale & $3J + 36 = 156$ & (con $J = 40$) \\
\bottomrule
\end{tabular}

\subsection{Quantità derivate}

\[
  b_r = \frac{\eta_r}{\psi_r} \qquad \text{(generated quantity in Stan)}
\]

Il coefficiente $b_r$ è il rapporto tra la sensibilità della slope al livello
del campo ($\eta_r$) e la dispersione degli effetti casuali ($\psi_r$).

\subsection{Modello completo in notazione probabilistica}

\[
\begin{aligned}
\alpha_r &\sim \N(0,\;1) \\
\psi_r,\,\sigma_r &\sim \HN(0,\;1) \\
\eta_r &\sim \N(0,\;1) \\
\gamma_{r,k} &\sim \N(0,\;1) \quad k = 1,\ldots,K_W \\
\beta_{r,k} &\sim \N(0,\;1) \quad k = 1,\ldots,K_B \\
z_{\nu,r,j} &\sim \N(0,\;1) \quad j = 1,\ldots,J \\[4pt]
\mu_{r,ij} &= \alpha_r + z_{\nu,r,j}\,(\psi_r + \eta_r\,x_{\text{logB},i})
             + \vec{\gamma}_r^\top \vec{x}_{W,i}
             + \vec{\beta}_r^\top \vec{x}_{B,j} \\[4pt]
y_{r,ij} &\sim \N(\mu_{r,ij},\;\sigma_r^2)
\end{aligned}
\]

% =============================================================================
\section{Schema covariate}

\begin{center}
\begin{tabular}{@{}p{5.5cm}p{6.5cm}@{}}
\toprule
\textbf{Covariati within-field} $X_W$ & \textbf{Covariati between-field} $X_B$ \\
variano per osservazione & un valore per campo \\
\midrule
logBottom (profondità)\newline
\phantom{logBottom}→ anche nella slope prop. & OnFarm \\
Texture1 (ILR$_1$: sand/silt/clay) & Irrigate \\
Texture2 (ILR$_2$: silt/clay) & Fertilised \\
BulkDensity & N\_Natural \\
\bottomrule
\end{tabular}
\end{center}

% =============================================================================
\section{Grafo diretto aciclico (DAG)}

\begin{center}
\begin{tikzpicture}[
    node distance = 1.4cm and 2.0cm,
    lat/.style    = {circle, draw=myorange, fill=mylightorange,
                     thick, minimum size=1.0cm, font=\small},
    resp/.style   = {circle, draw=mygreen, fill=mylightgreen,
                     thick, minimum size=1.0cm, font=\small},
    hp/.style     = {rectangle, rounded corners=3pt,
                     draw=mygray, fill=white,
                     thick, minimum width=1.0cm, minimum height=0.7cm,
                     font=\small\itshape},
    fixed/.style  = {rectangle, draw=myblue, fill=mylightblue,
                     thick, minimum width=1.1cm, minimum height=0.7cm,
                     font=\small},
    >=Stealth
  ]

  % Hyperparametri
  \node[hp] (psi)   at (0, 0)    {$\psi_r$};
  \node[hp] (eta)   at (2, 0)    {$\eta_r$};
  \node[hp] (alpha) at (-2, 0)   {$\alpha_r$};
  \node[hp] (sigma) at (6, 0)    {$\sigma_r$};

  % Effetto casuale
  \node[lat] (znu) at (1, -2)    {$z_{\nu,r,j}$};

  % Predittore e risposta
  \node[lat] (mu)  at (3, -2)    {$\mu_{r,ij}$};
  \node[resp] (y)  at (5, -2)    {$y_{r,ij}$};

  % Covariati fissi
  \node[fixed] (XW) at (3, -4)   {$X_W$};
  \node[fixed] (XB) at (1, -4)   {$X_B$};
  \node[fixed] (gam) at (4.2, -3.2) {\small$\vec{\gamma}_r$};
  \node[fixed] (bet) at (2.2, -3.2) {\small$\vec{\beta}_r$};

  % Frecce
  \draw[->] (psi)   -- (znu);
  \draw[->] (eta)   -- (mu);
  \draw[->] (psi)   -- (mu);
  \draw[->] (znu)   -- (mu);
  \draw[->] (alpha) -- (mu);
  \draw[->] (mu)    -- (y);
  \draw[->] (sigma) -- (y);
  \draw[->] (XW)    -- (mu);
  \draw[->] (XB)    -- (mu);
  \draw[->] (gam)   -- (mu);
  \draw[->] (bet)   -- (mu);

  % Plate per j
  \begin{scope}[on background layer]
    \node[draw=myorange, dashed, rounded corners=6pt,
          fit=(znu), inner sep=8pt,
          label={[font=\footnotesize\color{myorange}]below right:$j=1\ldots J$}
         ] {};
  \end{scope}

  % Plate per ij
  \begin{scope}[on background layer]
    \node[draw=mygreen, dashed, rounded corners=6pt,
          fit=(mu)(y), inner sep=8pt,
          label={[font=\footnotesize\color{mygreen}]below right:$i=1\ldots N_j$}
         ] {};
  \end{scope}

\end{tikzpicture}
\end{center}

\medskip
\noindent
\begin{tabular}{@{}ll@{}}
\toprule
Colore & Significato \\
\midrule
Arancione (cerchio) & variabile latente / effetto casuale \\
Verde (cerchio) & risposta osservata \\
Grigio (rettangolo) & iperparametro (prior) \\
Blu (rettangolo) & covariato fisso \\
\bottomrule
\end{tabular}

% =============================================================================
\section{Configurazione MCMC}

\begin{tabular}{@{}ll@{}}
\toprule
Impostazione & Valore \\
\midrule
Software & Stan via \texttt{cmdstanr} \\
Catene & 4 \\
Iterazioni warmup & 2\,000 \\
Iterazioni sampling & 5\,000 \\
\texttt{adapt\_delta} & 0.97 \\
\texttt{max\_treedepth} & 11 \\
Seed & 2024 \\
Parametri totali & 156 (con $J = 40$) \\
\bottomrule
\end{tabular}

% =============================================================================
\section{Risultati}

\subsection{Parametro $\eta_r$: correlazione slope-intercetta}

\begin{center}
\begin{tabular}{@{}lrrll@{}}
\toprule
Risposta & $\hat{\eta}_r$ (mediana) & SD & CI 90\% & Interpretazione \\
\midrule
logSOC & $+0.209$ & $0.055$ & $[0.125,\; 0.302]$ & forte, ben identificato \\
logN   & $-0.051$ & $0.064$ & $[\approx{-}0.15,\; +0.06]$ & trascurabile \\
logP   & $+0.096$ & $0.049$ & $[0.019,\; 0.179]$ & debole ma presente \\
\bottomrule
\end{tabular}
\end{center}

\medskip
Il coefficiente $b_{\SOC} = \eta_{\SOC}/\psi_{\SOC} \approx 0.47$:
per ogni unità di deviazione standard nel livello del campo,
la slope di decadimento verticale cambia di circa $0.47\,\psi_{\SOC}$.

\subsection{Confronto modelli (LOO-CV)}

\begin{center}
\begin{tabular}{@{}llll@{}}
\toprule
Modello & Descrizione & $\Delta$ELPD & Ranking \\
\midrule
\textbf{M-SP} & RI + slope proporzionale & --- & \textbf{1°} \\
M-RI & Random intercept i.i.d. & $<0$ & 2° \\
M-GPS & GP spaziale + slope fissa & $<0$ & 3° \\
M-GP & GP spaziale completo & $<0$ & 4° \\
\bottomrule
\end{tabular}
\end{center}

\noindent La struttura proporzionale cattura il segnale principale con il minor numero di
parametri. Il processo Gaussiano spaziale non è necessario: l'eterogeneità between-field
è già spiegata dal random intercept.

\subsection{Selezione variabili (projpred)}

Reference model: M-SP. Metodo: forward search con PSIS-LOO interno.
logBottom sempre incluso (nell'intercetta casuale proporzionale).

I risultati dettagliati si trovano in \texttt{scripts/06\_variable\_selection.R}.
In generale, pochi covariati fissi sono necessari oltre alla struttura random:
la maggior parte della variabilità è catturata da $z_{\nu,r,j}$.

% =============================================================================
\section{File del progetto}

\begin{tabular}{@{}ll@{}}
\toprule
File & Contenuto \\
\midrule
\texttt{scripts/01\_preprocessing.R} & crop.csv → dati.rds (ILR, pulizia) \\
\texttt{scripts/01b\_eda\_texture.R} & EDA composizionale tessitura \\
\texttt{stan/m4rr\_v2\_ri\_slope\_mu.stan} & codice Stan M-SP \\
\texttt{stan/fit\_v2\_ri\_slope\_mu.rds} & fit MCMC salvato \\
\texttt{scripts/03\_model\_proportional\_slope.R} & fit, diagnostica, LOO \\
\texttt{scripts/04\_eda\_depth\_profiles.R} & motivazione EDA \\
\texttt{scripts/06\_variable\_selection.R} & selezione variabili \\
\texttt{CLAUDE.md} & riferimento tecnico per agenti AI \\
\texttt{docs/note\_progetto.md} & note per umani \\
\bottomrule
\end{tabular}

\end{document}
